\documentclass[10pt]{article}
\usepackage[margin=0.72in]{geometry}
\usepackage[T1]{fontenc}
\usepackage{lmodern}
\usepackage{graphicx}
\usepackage{booktabs}
\usepackage{hyperref}
\hypersetup{pdftitle={Jet completion: finite control and proposed next search},pdfauthor={Research follow-up for PR 2969}}
\title{Jet completion: finite control and proposed next search}
\author{Research follow-up for PR 2969}
\date{2026-11-01}
\begin{document}
\maketitle
\section*{Evidence scope}
The existing native exact-rational control at $L=4,8,16,32$ checks the original cardinal-polynomial kernel against its ordinary first-jet flattening control. Its recorded first-coefficient valuations differ by two at each size, while both normalized residue ranks stay 2. The four rows contain 5,440 nodal value/first-derivative checks in total. These are finite exact observations, not a runtime comparison or a uniform theorem.
\section*{Recorded finite observations}
\begin{center}
\begin{tabular}{rrrrr}\toprule
$L$ & Nodal checks & Original $v_2$ & Flattened $v_2$ & Both residue ranks\\\midrule
4 & 64 & 3 & 5 & 2\\
8 & 256 & 4 & 6 & 2\\
16 & 1,024 & 5 & 7 & 2\\
32 & 4,096 & 6 & 8 & 2\\\bottomrule
\end{tabular}
\end{center}
The values are copied from \texttt{native/expected.txt} and match \texttt{native/replay.log}. The prospective protocol identifies the coefficients: the original is linear in the scaled chart and the flattened control is quadratic. Flattening roughly doubles polynomial degree and adds two valuation bits in these cases; the data do not establish that this trade is beneficial after degree and precision are priced.
\begin{center}
\includegraphics[width=0.96\linewidth]{valuation.pdf}
\end{center}
\noindent\textbf{Figure 1.} Exact recorded $v_2$ by $L$ from the native control. Each point is one size, not a timing sample. Normalized residue rank is 2 for both series at every size. The source SVG is \texttt{valuation.svg}.
\section*{Proposals, separately}
The study proposes (1) replacing the scalar origin channel with a nonscalar matrix $R(T)$ whose off-node values vanish and whose crossing jets have low weighted rank, and (2) allowing free jets through $W_{new}(T)=W_{base}(T)+C(T)H(T)$. Neither construction has a supplied integral factorization certificate or a measured cost here. The finite flattening control is not evidence that either proposal works.

The next discriminating step is a bounded search that returns the candidate polynomial and entry-by-entry integral $U,V$ witnesses at every retained node and stipulated precision. Its candidate domain is $\deg H\leq 2L-1$; initial sizes are $L=4,8$, precisions $8,12,16$, and controls are $H=0$, the recorded flattening construction, and rejection of an invalid even-denominator inverse. Before searching, freeze rational coefficient and denominator budgets and price polynomial degree and precision. A successful small case would still need a uniform construction and a complete recursive proof.

The linked backlog card \texttt{XP-20261101-rank-directed-jet-completion} is draft, unfunded, and not scheduled. The runner, frozen protocol, resource estimate, and authorization to execute do not exist yet.
\section*{Sources and limits}
\begin{itemize}
\item Study README: proposals, control interpretation, formal scope, and deferred search inputs.
\item \texttt{native/PROTOCOL.txt}: four-size exact-rational control protocol; no timing run.
\item \texttt{native/expected.txt} and \texttt{native/replay.log}: recorded values and replay.
\item \texttt{lean/PROOF\_MAP.md}: formalized own-node value and first-derivative identities and their limits.
\end{itemize}
No new experiment, Lean build, ECDLP instance, or solver execution was performed for this report.
\end{document}
